\section{Conditional Dirichlet covariance}\label{sec:conditional-covariance}

The response argument in Section~\ref{sec:cell} identified the Gaussian
covariance for ordinary fixed-boundary domains. We now identify the
conditional covariance when internal oriented strands have also been
observed. A common sign wall does not make the two coupled heights equal
there. Instead, the squared common-wall errors in the comparison estimate,
together with thickness of the observed set, give zero trace for the
limiting response. The preceding analytic section then identifies the
response as the Dirichlet Green solution, and exponential tilting gives
the covariance.
This argument uses no interface convergence or height-gap theorem.

\subsection{The uniform input and the conditioning class}

For each mesh $n^{-1}$ let $\mathcal S_n$ be a family of deterministic
oriented sign-string configurations with ordinary one-sided sign walls,
fixed Dirichlet values in $[-1/2,1/2]$, and positive-volume conditional
height polytopes.  The sign strings are of the type specified in
\ref{f:conditioned-sparse}; no further exact-height transcript is fixed.
Denote their conditional
laws by $P_{n,s}$, $s\in\mathcal S_n$.  Realize each observed set as a
closed geometric set $K_{n,s}\subset\overline Q$ containing $\partial Q$:
join the observed lattice vertices along their strings, include the
exterior fixed region, and use a fixed $O(n^{-1})$ convention at faces.
Such conventions have the same Hausdorff limits.  Require the common
mesh thickness constants of Definition~\ref{t:thickness}.  Write
$V_{n,s}=Q\setminus K_{n,s}$.

The following are the exact forms of the preceding estimates used here.
They make the scope of the uniform assertion explicit.  All constants may
depend on a fixed compact family $\mathcal C\subset C_c^\infty(Q)$ of
real test functions, but not on $n$ or $s$.
\begin{enumerate}
\item\label{t:input-moments}
The centered Riemann pairings
\[
 X_{n,s}(f)=\frac1{vn^2}\sum_x f(x)
                 \bigl(h_n(x)-\E_{n,s}h_n(x)\bigr)
\]
obey $\sup_{n,s,f\in\mathcal C}\E_{n,s}e^{a|X_{n,s}(f)|}<\infty$
for some $a>0$.  The estimate also holds for the fixed finite linear
combinations used for joint tests.  This is the form of
\ref{f:averages} used below.
\item\label{t:input-extension}
For each fixed $|t|<a/2$ couple $P_{n,s}$ to its tilt by
$e^{tX_{n,s}(f)}$ using the common sign walls.  For their difference
$z_{n,s,t}$, the filled extension $U_{n,s,t}$ supplied by
\ref{f:hole} and \ref{c:comparison} satisfies every hypothesis of
Proposition~\ref{t:varying-trace}, including
\eqref{t:original-filled}, uniformly in $s$ and in $f\in\mathcal C$.
In particular its energy includes squared common-wall errors.  Its
inverse-weight spatial moments have deterministic bounds off exceptional
events whose contributions vanish.  Expectation-only bounds for a random
weight constant are not silently multiplied by an energy known only in
$L^1$.
\item\label{t:input-flux}
For every deterministic sequence $s_n,f_n$ with $f_n$ converging in
$\mathcal C$ and $K_{n,s_n}$ converging in Hausdorff distance to $K$, the
bulk identity~\ref{b:cell} identifies each limit of the mean response by
\begin{equation}\label{t:weak-response}
 A\int_{Q\setminus K}\nabla m\cdot\nabla\phi\dd x
   =\frac{t}{v}\int_{Q\setminus K}f\phi\dd x,
 \qquad \phi\in C_c^\infty(Q\setminus K),
\end{equation}
with the same scalar $A$, positive by~\ref{b:response}.  The pairings with the actual fields and their
piecewise affine interpolants differ by $o(1)$ in the needed expectations.
\end{enumerate}

These clauses are used below as conclusions of the indicated finite-volume
results, not as consequences of the trace theorem.  In particular,
interior homogenization by itself supplies neither
\ref{t:input-extension} nor uniformity over $\mathcal S_n$.

We explain the retained-site requirement in
Proposition~\ref{t:varying-trace}. Outside the exceptional event, every
filled hole has diameter at most a fixed power of $\log n$ in lattice
units. Mesh thickness ensures that an observed component through a
vertex in such a hole reaches a fixed positive physical scale. It
therefore has an observed graph path exiting the hole, whose first
good rim vertex remains observed and lies within a polylogarithmic
lattice distance. Move the testing mass to that vertex. At a common
fixed pin its difference is zero; at a retained common sign wall its
squared value appears in \eqref{c:energy}. The rim construction retains
these values and has the stated reciprocal-weight moments. Thus this
movement verifies the anchor hypothesis, including for internal
components; thickness rules out an entire component hidden inside a
microscopic filled hole.

\begin{theorem}[Conditional covariance for deterministic strings]
\label{t:covariance}
Suppose the class $\mathcal S_n$ and the estimates
\ref{t:input-moments}--\ref{t:input-flux} hold, and set
$\sigma^2=(vA)^{-1}$.  Along every deterministic sequence
$s_n\in\mathcal S_n$ with $K_{n,s_n}\to K$ in Hausdorff distance,
\begin{equation}\label{t:conditional-mgf}
 \log\E_{n,s_n}e^{tX_{n,s_n}(f)}
 \longrightarrow
 \frac{t^2\sigma^2}{2}\langle f,G_{Q\setminus K}f\rangle
\end{equation}
for $t$ in a neighborhood of zero.  The conclusion holds for convergent
sequences of tests from a compact family and jointly for finitely many
tests.  In particular,
\begin{equation}\label{t:uniform-covariance}
 \sup_{s\in\mathcal S_n}\sup_{f,g\in\mathcal C}
 \left|\operatorname{Cov}_{n,s}(X_{n,s}(f),X_{n,s}(g))
       -\sigma^2\langle f,G_{V_{n,s}}g\rangle\right|
 \longrightarrow0.
\end{equation}
There is no requirement that $K$ have zero area or that the test supports
be disjoint from $K$.
\end{theorem}
\begin{proof}
Fix a deterministic sequence as in the statement.  For fixed $t$,
\ref{t:input-extension} and Proposition~\ref{t:varying-trace} give
subsequential compactness of the mean filled response in $L^p$, for a
fixed $3/2<p<2$, and a zero $1/2$-Frostman trace on $K$.  The original
response has the same distributional limit.  By
\eqref{t:weak-response} and Theorem~\ref{t:trace}, that limit is uniquely
\begin{equation}\label{t:response}
 m_t=\frac{t}{vA}G_{Q\setminus K}f.
\end{equation}
It is zero almost everywhere on $K$, including its positive-area parts.
This proves convergence against general smooth $f$, rather than only
against tests separated from the strings.  Uniqueness removes the need
to choose a further subsequence.

Set $L_n(t)=\log\E_{n,s_n}e^{tX_{n,s_n}(f)}$.  Differentiation of the
finite-volume integral gives
\[
 L_n'(t)
  =\frac1{vn^2}\sum_xf(x)
     \bigl(\E_{n,s_n,t}h_n(x)-\E_{n,s_n}h_n(x)\bigr)
  \longrightarrow t\sigma^2\langle f,G_{Q\setminus K}f\rangle.
\]
The factor $1/v$ in the source is necessary: the tilt adds the microscopic
drift $t f(x)/(vn^2)$ to the reflected bath.  No boundary term is added,
since the walls and the exterior values are common to both laws.

The exponential moment in \ref{t:input-moments} makes these derivatives
uniformly bounded for $|t|\le a/4$.  Indeed the denominator of a centered
moment generating function is at least one, and
\[
 |L_n'(t)|\le
 \E_{n,s_n}\bigl[|X_{n,s_n}(f)|e^{a|X_{n,s_n}(f)|/4}\bigr]\le C_a.
\]
Integrating from zero proves \eqref{t:conditional-mgf}.  The same
exponential bound gives tightness and uniform integrability of every
fixed moment.  Each limiting law is the Gaussian determined by this
moment generating function near zero; hence the second moments converge.
Applying this to linear combinations and polarizing gives cross
covariances and the finite-dimensional conclusion.

The proof also works with $f_n$ in a compact family converging to $f$.
Choose $d_n\downarrow0$ sufficiently slowly that
$(f_n-f)/d_n$ belongs to a compact family in $C_c^\infty(Q)$.  Apply
\ref{t:input-moments} to this additional compact family.  Linearity then
makes the change of the pairing tend to zero in every fixed moment.
This uses the stated estimates for each fixed compact family, with its
own constants, and requires no uniform constant over all such families.
To obtain
\eqref{t:uniform-covariance}, suppose its supremum did not tend to zero
and select a deterministic sequence realizing a fixed positive error.
Hausdorff compactness and compactness of the test family give a convergent
subsequence.  The convergence already proved contradicts that error,
once the continuity of the Green form is checked.  That continuity is
proved in the following lemma and is purely analytic.
\end{proof}

\begin{lemma}[Green continuity for this class]\label{t:green-continuity}
If compact sets $K_n$ containing $\partial Q$ satisfy the common thickness
condition and tend in Hausdorff distance to $K$, then
\[
 \langle f_n,G_{Q\setminus K_n}g_n\rangle
 \longrightarrow\langle f,G_{Q\setminus K}g\rangle
\]
whenever bounded smooth tests $f_n,g_n$ converge uniformly to $f,g$.
The same assertion holds for mesh-thick approximants.
\end{lemma}
\begin{proof}
Put $v_n=G_{Q\setminus K_n}g_n$ and extend by zero.  Their variational
identities and the Poincare inequality on $Q$ give a common $H^1$ bound.
After taking a subsequence they converge strongly in $L^2$ and weakly in
$H^1$ to $v$.  Every compact subset of $Q\setminus K$ is eventually in
$Q\setminus K_n$, so $-\Delta v=g$ there.

The zero traces pass to $K$ by the deterministic version of the proof of
Proposition~\ref{t:varying-trace}: use the continuous, uniformly Frostman
approximating measures from Lemma~\ref{t:transport}, before its optional
projection onto lattice vertices.  The $H^1_0(Q\setminus K_n)$
approximation by compactly supported smooth functions and
Lemma~\ref{t:frostman} show that $v_n$ has zero trace on these measures.
Atomic vertex traces of $v_n$ are neither asserted nor used.
There are no wall errors or filling maps in this case.  The $H^1$ bound
implies the required $W^{1,p}$ bound, and the averaging estimate passes
the traces.  Theorem~\ref{t:trace} identifies $v=G_{Q\setminus K}g$ on
all of $Q$.  Thus every subsequence has this same limit.  Strong $L^2$
convergence and uniform convergence of $f_n$ prove the assertion.
For mesh-thick sets perform the measure construction only at scales
larger than $c/n$; its errors tend to zero before the fixed averaging
scale is removed.
\end{proof}

\subsection{Random strings and boundary sampling}

\begin{corollary}[Transfer through an exact conditional kernel]
\label{t:random-covariance}
Let $S_n$ take values in $\mathcal S_n$, and suppose
\begin{equation}\label{t:kernel}
 \operatorname{Law}(h_n\mid S_n=s)=P_{n,s}
 \quad\hbox{for almost every }s.
\end{equation}
Then, for fixed tests or compact families as in Theorem~\ref{t:covariance},
the conditional covariance differs in $L^1$ from
$\sigma^2\langle f,G_{V_{n,S_n}}g\rangle$ by $o(1)$.
If $K_{n,S_n}$ converges in distribution to $K$, these random Green forms
converge in distribution and in expectation to those on $Q\setminus K$.
Convergence in $L^1$ to a limiting random form additionally requires that
the sets converge in probability on a common probability space.
\end{corollary}
\begin{proof}
Condition on $S_n$ and apply the deterministic uniform error
\eqref{t:uniform-covariance}; its expectation is no larger than that
error.  Lemma~\ref{t:green-continuity} gives convergence of the random
Green forms.  They are uniformly bounded for bounded tests, because the
Dirichlet Green quadratic form of every subdomain is bounded above by
that of $Q$.  This also gives convergence of expectations and the stated
$L^1$ assertion under a common coupling.
\end{proof}

For an exploration that queries signs and makes every next query or
stopping decision from its previous answers and independent randomness,
the full transcript has the conditional-kernel property: its decision
indicators introduce no factor depending on an unqueried height.
Disintegrating the finite-dimensional Lebesgue density therefore leaves
exactly the uniform density with the recorded signs imposed.  An
exact-value exploration gives instead the corresponding frozen-value
kernel; it is covered only when the finite-volume estimates have been
proved for that class.  An arbitrary field-dependent selection of a
string need not satisfy \eqref{t:kernel}.  For example a label chosen by
the sign of a standard Gaussian variable changes its conditional
variance from one to $1-2/\pi$, although either deterministic label could
have been assigned the original reference kernel.  Uniform estimates for
reference kernels do not identify the law of such a selection.

\begin{proposition}[Vanishing covariance near a continuum]
\label{t:boundary-variance}
Let $K$ satisfy Definition~\ref{t:thickness} and let $V=Q\setminus K$.
Let $\mu_\eps$ be positive measures of mass at most one, supported in $V$
within distance $C\eps$ of $K$, with a common $1/2$-Frostman constant.
Then
\begin{equation}\label{t:boundary-energy}
 \langle\mu_\eps,G_V\mu_\eps\rangle\longrightarrow0.
\end{equation}
The conclusion holds for any unnormalized restriction of such a measure.
It holds uniformly after adding further observed sets: if $K'\supset K$
and $V'=Q\setminus K'$, the Green quadratic form on $V'$ is bounded above
by that on $V$.
\end{proposition}
\begin{proof}
The Green function of $V$ is bounded above by that of a fixed disk
containing $Q$, hence by $C(1+\log^+(1/|x-y|))$.  Dyadic summation with
the Frostman bound gives
\[
 \iint_{|x-y|<\rho}G_V(x,y)\dd\mu_\eps(x)\dd\mu_\eps(y)
 \le C\rho^{1/2}(1+|\log\rho|).
\]
For $|x-y|\ge\rho$, a Brownian path started at $x$ must leave a
$c\rho$-neighborhood of a nearest point of $K$ before it contributes to
the Green function at $y$.  Beurling bounds that probability by
$C(\eps/\rho)^{1/2}$.  On the exit circle the Green function is at most
$C(1+|\log\rho|)$.  Thus the remaining contribution is at most
$C(\eps/\rho)^{1/2}(1+|\log\rho|)$.  First send $\eps$ to zero and
then $\rho$ to zero.  Restricting the measure preserves all these bounds.

For the last assertion, the inclusion $H^1_0(V')\subset H^1_0(V)$ and
the variational formula
\[
 \langle f,G_Vf\rangle
 =\sup_{u\in H^1_0(V)}
       \left\{2\langle f,u\rangle-\int_V|\nabla u|^2\right\}
\]
give the quadratic-form inequality, also for signed tests.  Apply it
first to smooth measures and then use the finite logarithmic energy to
pass to the stated measures.
\end{proof}

\begin{corollary}[Local boundary variance, independently of ambient size]
\label{t:local-boundary-variance}
Fix nested squares $Q_0\Subset Q_1\Subset Q_2\Subset Q_3$ and positive
geometric separation constants.  Consider finite triangular-lattice
domains of mesh $n^{-1}$, with arbitrary ambient diameter, fixed data in
$[-1/2,1/2]$, and deterministic sign strings.  Suppose a connected
observed piece of diameter at least $c>0$ lies within bounded distance of
$Q_2$.  Assume the local hypotheses of~\ref{f:averages} and
\ref{c:comparison}, with constants independent of ambient diameter.
Assume also that restricting the observed system to $Q_3$ and adjoining
its artificial exterior gives the common mesh thickness constants of
Definition~\ref{t:thickness}.  This holds for connected strands attached
to the original exterior, or for components of uniformly positive
diameter: cut pieces meeting the artificial boundary join that boundary,
and pieces entirely inside are whole original components.

For each fixed $\eta>0$, let $\mu_{n,\eta}$ be a deterministic smooth
probability test supported in $Q_0$, at distance at most $C\eta$ from
the observed set in $Q_1$.  Suppose these measures have a common
$1/2$-Frostman constant, independent of $n,\eta$, and for fixed $\eta$
their densities range over a compact smooth-test family.  If $X_n(\mu)$
denotes the Riemann pairing of the centered field with that density, then
\begin{equation}\label{t:local-variance-limit}
 \lim_{\eta\downarrow0}\limsup_{n\to\infty}
 \sup\operatorname{Var}X_n(\mu_{n,\eta})=0.
\end{equation}
The supremum is over the stated local geometries, tests, and arbitrary
ambient domains.  The inner limit is always taken at fixed $\eta$.
\end{corollary}
\begin{proof}
Construct an auxiliary height law in the portion of the original domain
inside $Q_3$.  Retain all original fixed data and all sign constraints
on its free vertices, and put zero values on the new artificial
boundary.  At a vertex already fixed by the original problem retain
that original value.  Drop sign walls on vertices newly fixed at the
artificial rim; their values are now prescribed to be zero.  The two
models have exactly the same graph,
fixed data, and sign walls in $Q_2$ once $n$ is large.  A strict feasible
filling for the auxiliary free coordinates is obtained by assigning
small signed values at the sign sites and small values at the other
free sites: each original fixed value has absolute value at most $1/2$.
Thus the auxiliary conditional polytope has positive volume.  At a cut
string endpoint, a free sign site whose opposite neighbor was removed
is adjacent to the zero artificial boundary; it therefore remains
bounded.  Every truncated observed component reaching the artificial
boundary joins the exterior observed set.  The stated thickness
assumption consequently puts this auxiliary law within the bounded
domain class of Theorem~\ref{t:covariance}.

Let $h_n^{\rm loc}$ denote the auxiliary field.  Put it and the original
field $h_n$ in the stationary synchronous joining supplied by
\ref{c:comparison}, using the common walls in $Q_2$.  The artificial
outer data may differ, as permitted in that local comparison theorem.
The original marginal has uniform local average bounds from the nearby
observed piece.  The auxiliary marginal has them from its artificial
boundary, or the retained observed piece.  If $U_n$ is the filled extension of
$z_n=h_n-h_n^{\rm loc}$ in $Q_1$, its energy and squared Sobolev norm
satisfy
\begin{equation}\label{t:local-difference-energy}
 \E E_n+\E\|U_n\|_{W^{1,p}(Q_1)}^2\le C_p,
 \qquad 3/2<p<2,
\end{equation}
with constants independent of the original ambient diameter and of
$\eta$.  The coupling is constructed before choosing the test.

Transport $\mu_{n,\eta}$ to the nearby common observed continuum,
then to its mesh vertices and past the holes to retained vertices.
The transport in Lemma~\ref{t:transport} moves points by $O(\eta)$;
the last two movements are $O(\ell_n/n)$ and give the coarse Frostman
bound~\eqref{t:coarse-frostman}.  The squared anchor estimate is
\[
 \left(\sum_x\nu_n(x)|U_n(x)|\right)^2
 \le E_n\sum_x\frac{\nu_n(x)^2}{w_x}
 \le CE_n n^{-1/2+3/(2q)}(\log n)^C,
 \qquad q>3.
\]
Its expectation tends to zero.  Compare each testing value with its
$\eta$ average as in Lemmas~\ref{t:frostman} and~\ref{t:transport}.
The bottom mesh scales contribute $o(1)$ at fixed $\eta$.  Squaring
these inequalities and using~\eqref{t:local-difference-energy} gives
\begin{equation}\label{t:local-difference-trace}
 \limsup_{n\to\infty}
 \E\left|\langle\mu_{n,\eta},U_n\rangle\right|^2
 \le C_p\eta^{2\beta},\qquad
 \beta=1-\frac{3}{2p}>0.
\end{equation}
All these inequalities are samplewise before expectation; the moved
measure may depend on the holes.  No independence between that measure
and $U_n$ is required.

It remains to transfer the upper bound to the original difference.
At fixed $\eta$, put
\[
 Z_n=\langle\mu_{n,\eta},h_n-h_n^{\rm loc}\rangle,
 \qquad W_n=\langle\mu_{n,\eta},U_n\rangle.
\]
The original--filled $L^1$ estimate in~\ref{c:comparison}, together
with the bounded test density at this fixed scale, gives
$\E|Z_n-W_n|\to0$.  The marginal fourth-moment estimates in
\ref{f:averages} make $Z_n^2$ uniformly integrable.  Their constants
may depend on $\eta$, which is harmless for this order of limits.
For $M>0$, the function $x\mapsto\min\{x^2,M\}$ is
$2\sqrt M$-Lipschitz, so
\[
 \limsup_n\E\min\{Z_n^2,M\}
 \le\limsup_n\E W_n^2\le C_p\eta^{2\beta}.
\]
Letting $M\to\infty$ proves
$\limsup_n\E Z_n^2\le C_p\eta^{2\beta}$.  No $L^2$ estimate for the
original--filled error has been assumed.

The auxiliary domain has uniformly bounded diameter in these scaled
coordinates.  At each fixed $\eta$, Theorem~\ref{t:covariance} applies
uniformly to its compact family of tests, and
Proposition~\ref{t:boundary-variance} makes its limiting tested variance
tend to zero as $\eta\downarrow0$.  Finally the coupling inequality
\[
 \sqrt{\operatorname{Var}\langle\mu_{n,\eta},h_n\rangle}
 \le
 \sqrt{\operatorname{Var}\langle\mu_{n,\eta},h_n^{\rm loc}\rangle}
       +\|Z_n\|_2
\]
and~\eqref{t:local-difference-trace} prove
\eqref{t:local-variance-limit}.  The proof identifies no Poisson
solution on an unbounded domain and imposes no condition at infinity.
\end{proof}

For the harmonic tests used later, first stop Brownian motion a positive
distance $\eps$ from the observed set and smear the exit law into balls
of radius $c\eps$, $c<1/2$, contained in the open component.
Lemma~\ref{t:exit} gives the required Frostman bound, which smearing
preserves.  At this fixed scale the smoothed densities form a compact
family of smooth tests with uniform derivative bounds.  Apply
Theorem~\ref{t:covariance}, including if further strings meet their
supports, and only afterwards send $\eps$ to zero using
Proposition~\ref{t:boundary-variance}.  The zero-response assertion on
positive-area limits is what permits a fixed test to meet the additional
closed set.  No estimate uniform in $\eps$ before taking the lattice
limit has been used.

\begin{remark}[Finite-volume dependencies and conditioning scope]
For the deterministic class specified above, clause~\ref{t:input-extension}
uses the following parts of the finite-volume proof.
Corollary~\ref{f:conditioned-sparse} supplies sparse initial contacts
uniformly over the sign strings.  The forcing and loss allocation in
Lemma~\ref{c:force} and equations~\eqref{c:force-loss}--\eqref{c:allocation}
retain the common-wall square penalties in~\eqref{c:energy}.
Equation~\eqref{c:conditional-moments} gives deterministic spatial
inverse-weight bounds outside arbitrarily rare exceptional events.
The interpolation and absorption estimates
\eqref{c:interpolation}--\eqref{c:absorb}, followed by the hole estimate
\ref{f:hole}, give the squared $W^{1,p}$ bound in~\eqref{c:compactness}
for every $1<p<2$; in particular one may choose $p>3/2$.
The all-edge estimate~\eqref{c:all-edge-l1} and the subsequent
fixed-construction hole truncation give the vanishing original--filled
$L^1$ error, not merely its boundedness.  Lemma~\ref{c:exterior-zero}
preserves the common zero exterior, while the common-wall penalties and
the thickness hypothesis give the internal traces through
Proposition~\ref{t:varying-trace}.  Lemma~\ref{f:tilts} and the small-drift
completion of the comparison proof preserve these estimates under the
small smooth tilts used here.

This argument concerns ordinary one-sided sign constraints and the stated
range of original fixed data.  It does not include shrinking two-sided
wall intervals or assert uniform estimates after an arbitrary exact-value
transcript is fixed.  The adapted exact-value comparisons in
Lemma~\ref{c:freezing} are annealed over their actual transcript; they do
not enlarge the deterministic class $\mathcal S_n$ of
Theorem~\ref{t:covariance}.  Random use of that theorem still requires the
exact conditional-law identity in Corollary~\ref{t:random-covariance}.
The normalization of the resulting covariance is the $(-\Delta)^{-1}$
convention used in~\cite{SScontinuum}.
\end{remark}
